One might want to jointly minimize the loss $\tilde{\mathcal{L}}^{(n)}_2$, which is the empirical approximation of $\mathcal{L}$. However, this fails to learn the true structural function $f_{\mathrm{struct}}$, as shown in this appendix.

Figure~\ref{fig:learning_fail} shows the learning curve obtained when we jointly minimize $\theta_X$ and $\theta_Z$ with respect to $\tilde{\mathcal{L}}^{(n)}_2$, where $\mathcal{L}_{\mathrm{test}}$ is the empirical approximation of
\begin{align}
    {L}_{\mathrm{test}} = \expect[X]{\|f_{\mathrm{struct}}(X) - (\hat{\vec{u}}^{(n)})^\top\vec{\psi}_{\theta_X}(X)\|^2}.
\end{align}

We observe in Figure~\ref{fig:learning_fail} that the decrease of stage 2 loss does not improve the performance of the learned structural function. This is because the model focuses on learning the relationship between instrument $Z$ and outcome $Y$, while ignoring treatment $X$. We can see this from the fact that stage 1 loss becomes large and unstable. This is against the goal of IV regression, which is to learn a causal relationship between $X$ and $Y$.

On the other hand, Figure~\ref{fig:learning_suc} shows the learning curve obtained using DFIV. Now, we can see that stage 2 loss matches $\mathcal{L}_{\mathrm{test}}$. Also, we can confirm that stage 1 loss stays small and stable.


\begin{figure}[t]
    \begin{minipage}{0.48\hsize}
    \centering
    \includegraphics[width=0.95\textwidth]{Appendix/learning_curve_fail.pdf}

    \caption{Learning curve of joint minimization}
    \label{fig:learning_fail}
    \end{minipage}
    \hfill
    \begin{minipage}{0.48\hsize}
        \centering
        \includegraphics[width=0.95\textwidth]{Appendix/learning_curve_suc.pdf}
        \caption{Learning curve of DFIV}
        \label{fig:learning_suc}
    \end{minipage}
\end{figure}

